% Pillar-7: population-scale U-shape, mastery/incapacity aliasing, PCD
\section{The U-Shape at Population Scale, and Why ZVF Cannot Steer}
\label{sec:p7-ushape}

\subsection{The 368-run audit, and what it can and cannot show}
\label{sec:p7-audit}
Theorem~1 is a statement about one prompt. For a population of prompts with
success probabilities $p_i$ and conditionally i.i.d.\ rollouts, it gives
$\E[\mathrm{ZVF}]=\E_i[h_G(p_i)]\ge h_G(\bar p)$ (convexity of $h_G$), so batch
ZVF is high whenever most prompts sit near $p_i\approx0$ \emph{or}
$p_i\approx1$, and its excess over $h_G(\bar p)$ measures the
\emph{dispersion} of prompt difficulty. On the 600 GSM8K groups (three seeds,
$G{=}8$), observed ZVF is $0.158$, $h_8(\bar p)=0.054$, and a beta-binomial
fit with intra-prompt correlation $0.16$ predicts $0.170$. ZVF is the
Advantage Collapse Rate of \citet{he2026advantage} up to aggregation
granularity, and the Jensen excess over the Bernoulli prediction was proved by
\citet{nie2026gradstarvation}; we claim neither, only the use of ZVF as a
sampler-control input.

A W\&B audit (project \texttt{zvf-audit}; 368 runs, 350 with a run-level
summary, 352 with per-step histories) spans \textbf{nine} model labels
(Table~\ref{tab:p7-ushape}). Sorted by \emph{final-step} training reward,
per-model mean ZVF is $0.95$--$0.97$ for Llama-3.2-1B/3B (reward $\approx0.01$)
and Qwen3-30B-A3B-Instruct (reward $0.95$), $0.25$--$0.29$ for Qwen3-32B/8B
(reward $0.35$--$0.50$), and intermediate for Nemotron-3-Nano-30B
($0.50$ at $0.77$), Llama-3.1-8B-Instruct ($0.65$ at $0.83$), Qwen3-4B-Instruct
($0.87$ at $0.86$), and gpt-oss-20b ($0.79$ at $0.90$). An earlier revision
showed seven of these rows and omitted Llama-3.1-8B-Instruct and gpt-oss-20b.

\begin{table}[t]
\centering
\small
\caption{Per-model ZVF in the 368-run \texttt{zvf-audit} corpus, sorted by
mean \emph{final-step} training reward. $n$ = runs with a summary;
$h_G(p_{\mathrm{step}})$ and excess are step-level means (excess $=$ step ZVF
$-\,h_G$(step reward)). Batch, $G$, and step budget differ by model (last
column), so the rows are \emph{not} a controlled comparison. Script:
\texttt{analysis\_publication\_fixes/zvf\_confound.py}.}
\label{tab:p7-ushape}
\begin{tabular}{lrccccl}
\toprule
Model & $n$ & Final reward & ZVF & $h_G(p_{\mathrm{step}})$ & Excess & Config (batch; steps) \\
\midrule
Llama-3.2-1B             & 101 & 0.01 & 0.97 & 0.97 & 0.00 & 8; 3--5 \\
Llama-3.2-3B             & 106 & 0.01 & 0.95 & 0.94 & 0.00 & 6--32; 2--12 \\
Qwen3-32B                &   6 & 0.35 & 0.25 & 0.10 & 0.12 & 6--32; 2--12 \\
Qwen3-8B                 &   6 & 0.50 & 0.29 & 0.09 & 0.19 & 6--32; 2--12 \\
Nemotron-3-Nano-30B      &   6 & 0.77 & 0.50 & 0.32 & 0.20 & 6--32; 2--12 \\
Llama-3.1-8B-Instruct    &   6 & 0.83 & 0.65 & 0.48 & 0.19 & 6--32; 2--12 \\
Qwen3-4B-Instruct-2507   & 107 & 0.86 & 0.87 & 0.47 & 0.41 & 6--32; 2--12 \\
gpt-oss-20b              &   6 & 0.90 & 0.79 & 0.49 & 0.32 & 6--32; 2--12 \\
Qwen3-30B-A3B-Instruct   &   6 & 0.95 & 0.96 & 0.64 & 0.31 & 6--32; 2--12 \\
\bottomrule
\end{tabular}
\end{table}

\paragraph{Confound analysis.} The nine rows are not a designed comparison.
The three small models contribute 314 of 350 runs, and 300 of those ran
5 steps at batch 8; 30 of the 36 runs of the other six models ran 12 steps at
batch 16 or 32, with $G$ varying within model. We therefore regressed ZVF on
model, batch, $G$, learning rate, and reward (linear $+$ quadratic), using
the pulled step histories
(\texttt{analysis\_publication\_fixes/zvf\_confound.py}; outputs in
\texttt{experiments/results/publication\_fixes/zvf\_confound/}).
At run level ($n=350$, HC3 SEs), a reward quadratic alone gives $R^2=0.42$;
model identity alone gives $0.87$; model $+$ configuration gives $0.90$; adding
the reward terms raises this by only $0.012$ (incremental $F=21.6$,
$p<10^{-8}$). At step level ($n=2{,}004$, SEs clustered by run), model $+$
batch $+$ $G$ gives $R^2=0.77$, adding reward terms $0.82$, and adding
$h_G(p_{\mathrm{step}})$ $0.82$. Within models, step ZVF correlates positively
with $h_G(p_{\mathrm{step}})$ in all nine (pooled within-model
$r=+0.49$), whereas its correlation with raw reward changes sign across models
($-0.89$ to $+0.70$). The cross-model U is thus nine points whose
position is set mostly by model identity and configuration. Because reward
is itself a property of the model, the regression cannot separate ``low ZVF
because interior difficulty'' from ``low ZVF because of the Qwen3-32B/8B
configuration''; the trough is not evidence of a reward effect.
The model-mean quadratic ($n=9$, $R^2=0.96$) should be read in that light.

\paragraph{What the U-shape still shows.} The \emph{endpoints} follow from the
identity: models at reward $\approx0$ sit on $h_G$ (excess $0.00$), and models
near reward $1$ have ZVF near $1$. The interior is dominated by excess over
$h_G(\bar p)$, i.e.\ difficulty dispersion, which is largest for
Qwen3-4B-Instruct ($0.41$). A controller that observes only a high ZVF still
cannot tell mastery (stop or harden the data) from incapacity (soften the data
or grow $G$), since the two demand \emph{opposite} interventions. This aliasing
follows from the symmetry $h_G(p)=h_G(1-p)$ and does not need the
population U to be causal: \emph{ZVF is a usable alarm but a poor steering
signal.}

\paragraph{Prospective ZVF prediction test (placeholder).}
% PLACEHOLDER -- ZVF PREDICTION EXPERIMENT (Colab agent running). Results will
% land in platform_hybrid/experiments/results/publication_fixes/. Replace this
% paragraph with: per-prompt p_i estimated on an independent sample, predicted
% batch ZVF = mean_i h_G(p_i) at G in {4,8,16}, observed ZVF, error and CI.
\textbf{[Placeholder: pre-registered held-out test of
$\widehat{\mathrm{ZVF}}=\frac1n\sum_i h_G(\hat p_i)$, results pending in
\texttt{experiments/results/publication\_fixes/}.]}

\subsection{PCD: the pairwise-contrast density}
\label{sec:p7-pcd}
The disambiguating statistic falls out of the same model. Define the
\emph{pairwise-contrast density} as the expected within-group reward
variance,
\begin{equation}
  \mathrm{PCD} \;=\; \E_x\bigl[\hat p_x(1 - \hat p_x)\bigr],
  \qquad
  \E[\mathrm{PCD}] \;=\; \frac{G-1}{G}\,\E_x\bigl[p_x(1-p_x)\bigr],
  \label{eq:p7-pcd}
\end{equation}
where the identity (Appendix~\ref{sec:p7-appendix}) follows from
$\E[\hat p(1-\hat p)] = (1 - \nicefrac{1}{G})\,p(1-p)$ for a binomial
sample. PCD measures the \emph{magnitude} of within-group contrast rather
than its mere existence, and it is proportional to the very $p(1-p)$ factor
that T3 ties to gradient mass: PCD is, up to the $\frac{G-1}{G}$ factor, the
average usable contrast per prompt. On the binned GSM8K reward tensors, PCD
traces the parabola the theory predicts (peak $0.25$ at $\hat p_x = 0.5$,
falling symmetrically to $0$ at the endpoints) while the ZVF indicator is a
flat $0$ everywhere in the interior and jumps to $1$ only at the endpoints
(\texttt{experiments/results/pcd\_vs\_zvf\_shape.tsv}): ZVF sees only the
cliff edges, PCD sees the slope.

\subsection{The micro-jitter falsification}
\label{sec:p7-jitter}
The sharpest separation between the two statistics is an intervention, not a
correlation. Because ZVF is defined by \emph{exact} reward ties, adding
i.i.d.\ jitter $\eps \sim \mathrm{U}(0, 10^{-4})$ to every reward makes ties
measure-zero and drives measured ZVF to $0$, while shifting the expected
within-group variance by only $\mathrm{Var}(\eps) \leq 10^{-9}$
(Appendix~\ref{sec:p7-appendix}). Running this on $600$ prompt-groups from
the GSM8K reward tensors (\texttt{scripts/pcd\_vs\_zvf.py}): batch ZVF
collapses $0.158 \to 0.000$, while PCD is unchanged at the reported
precision ($0.1538$ before and after; observed shift $< 10^{-6}$). Any
pipeline that emits a tiny dense sub-reward---a length penalty, a format
bonus---silently performs this jitter and reports ZVF $= 0$ over a batch
that is, for learning purposes, more than $15\%$ dead. A controller keyed to
raw ZVF would read such a batch as healthy; one keyed to PCD would not. This
is also our operational answer to reward-shaping ambiguity: shaped rewards
do not invalidate the starvation theory, they merely break the ZVF
\emph{estimator} of it, and PCD is the estimator that survives.

One caveat keeps us honest about what PCD is for. On $80$ runs with
per-run summaries, mean training reward predicts the held-out outcome far
better (Spearman $\rho = 0.95$) than mean ZVF does ($\rho = 0.56$ on this
subset). The companion Pillar-2 paper's earlier $\rho = 0.27$ on a 23-cell
matrix is withdrawn (17 finite cells, 9 of them simulation rows, mixed outcome
definitions) and is not used here. PCD is proposed
as a \emph{control-loop input}---a live measure of usable contrast---not as
an outcome predictor; the decisive predictive horse race needs per-group
reward tensors logged for all anchors, which only the GSM8K runs currently
carry (Section~\ref{sec:p7-limitations}).
